
Experimental design validates three core predictions through five complementary experiments: (P1) lightweight features achieve >80\% accuracy on held-out validation, (P2) features generalize to edge case architectures, and (P3) features exhibit scale invariance and strong inter-family separation.

\subsubsection{Experimental Questions}

\textbf{H-E1 (Existence).} Do normalization counts and parameter-mass ratio achieve >80\% 3-way classification accuracy on held-out TIMM models?

\textbf{H-M1 (Mechanism: Normalization Fingerprinting).} Do CNNs exclusively use BatchNorm while Transformers use LayerNorm, with violation rates $\leq 15$\% per class?

\textbf{H-M2 (Mechanism: Parameter Allocation).} Does parameter-mass ratio R exhibit exceptional inter-family separation (Cohen's d > 1.0) between CNN and Transformer families?

\textbf{H-M3 (Mechanism: Checkpoint-Only Feasibility).} Can features be extracted from checkpoint files in <10 minutes with 0 MB GPU usage?

\textbf{H-C1 (Condition: Edge Case Robustness).} Do features maintain $\geq 70$\% accuracy on non-standard architectures with $\leq 15$\% degradation from baseline?

\subsubsection{Dataset Construction}

60 pre-trained models were curated from TIMM model zoo (version 0.9.12), ensuring diversity across architecture families and model scales. Model selection included 24 CNN models (ResNet-{18,34,50,101,152}, VGG-{11,16,19}, DenseNet-{121,161,201}, EfficientNet-{B0,B1,B2,B3}, MobileNetV3-{Small,Large}, RegNet-{Y\_400MF,Y\_800MF,Y\_1\_6GF,Y\_3\_2GF}, ConvNeXt-{Tiny,Small,Base}), 24 Transformer models (ViT-{Tiny,Small,Base,Large}, DeiT-{Tiny,Small,Base}, Swin-{Tiny,Small,Base}, BEiT-{Base,Large}, CaiT-{XXS24,XS24,S24}, LeViT-{128,192,256}, PoolFormer-{S12,S24,S36,M36,M48}), and 12 Hybrid models. Models span $5\times$ parameter range within families to test scale invariance.

Stratified 70/30 split produced 42 training models (17 CNN, 17 Transformer, 8 Hybrid) and 18 validation models (7 CNN, 7 Transformer, 4 Hybrid). Random seed fixed at 42 for reproducibility.

Edge case validation set included 12 additional models: 1 NormFree network (NFNet), 3 SENet variants, 5 RegNet extreme scales, 2 ViT extreme variants, and 1 unknown architecture. These violate standard normalization conventions or operate at extreme parameter scales.

Ground truth labels were derived from TIMM naming conventions with structural validation on 10-model sample. TIMM naming alignment was 40\% (violating initial $\geq 90$\% assumption), but this validates that features extract structural information independent of naming.

\subsubsection{Metrics}

For H-E1, macro-averaged accuracy (class-balanced average) is reported to avoid bias toward majority classes. Per-class precision and recall must meet $\geq 75$\% threshold. For H-M2, Cohen's d effect size quantifies inter-family separation strength (d > 0.8 = large effect). For H-M1, violation rates measure percentage of models not using dominant normalization type per class. For H-M3, extraction time per model and peak GPU memory usage are measured.

For scale invariance (H-E1), coefficient of variation (CV = std/mean) is computed across ResNet-{18,34,50,101,152} family, requiring CV < 0.15. For inter-family separation (H-M2), Welch's t-test for unequal variances is applied with Bonferroni correction for multiple comparisons. For edge case degradation (H-C1), accuracy drop from baseline is computed, requiring $\leq 15$\% relative degradation.

\subsubsection{Implementation}

All experiments use Python 3.9, PyTorch 2.1.0, scikit-learn 1.3.0, and TIMM 0.9.12. Checkpoints are cached locally. Feature extraction is sequential, taking ~61 seconds for 60 models. Training and evaluation complete in <2 minutes. Total experimental runtime: ~3 hours, dominated by one-time checkpoint downloads. No GPU required---all operations are CPU-only.
